% !TEX root = ../main.tex
% !TEX program = lualatex
\documentclass[../main.tex]{subfiles}
\IfSubfilesClassLoaded{\externaldocument{\subfix{../build/main}}}{}
% ====================
% File: 16_Milestones.tex
% ====================
\begin{document}
\IfSubfilesClassLoaded{\chapter{EDO Study Guide}}{}
% === MILESTONES (EDO 3.3.2) ===
\section{Acquisition Milestones and Phases}

\subsection{Summary}
\begin{description}
  \item[\textbf{DAS structure.}] The \ac{das} uses Milestones~A/B/C, five phases, and decision points (\ac{mdd}, \ac{cddv}, \ac{drfprd}, \ac{frp}) to pace major capability acquisitions~\autocite{edo-3-3-2-acq-milestones-phases-2025}.
  \item[\textbf{Baseline discipline.}] The \ac{apb} sets cost, schedule, and performance thresholds/objectives that define the program trade space and trigger deviation actions~\autocite{edo-3-3-2-acq-milestones-phases-2025}.
  \item[\textbf{Pre-systems acquisition.}] The \ac{mdd}, \ac{msa}, and \ac{aoa} converge on a validated materiel solution, initial acquisition strategy, and draft requirements~\autocite{edo-3-3-2-acq-milestones-phases-2025}.
  \item[\textbf{Systems acquisition.}] \ac{tmrr} and \ac{emd} reduce risk and mature designs before \ac{pd.phase}; \ac{ms}~B/C decisions commit the program to development and production~\autocite{edo-3-3-2-acq-milestones-phases-2025}.
  \item[\textbf{Navy overlays.}] Ships, \ac{scn} timing, and the Navy 2-pass/7-gate review shape milestone execution and resourcing expectations~\autocite{edo-3-3-2-acq-milestones-phases-2025}.
\end{description}

\subsection{Introduction: DAS and Urgent Needs}
The \ac{das} supports the National Defense Strategy through the \ac{aaf}, tailoring acquisition to program urgency and risk~\autocite{edo-3-3-2-acq-milestones-phases-2025}.
\begin{description}
  \item[\textbf{UON.}] \ac{uon} addresses capability gaps impacting ongoing or anticipated contingency operations; when validated by a single \ac{dod} Component it is a Component \ac{uon}~\autocite{edo-3-3-2-acq-milestones-phases-2025}.
  \item[\textbf{JEON.}] \ac{jeon} is a joint requirement identified by \ac{ccmd}, \ac{cjcs}, or \ac{vcjcs} for an anticipated contingency~\autocite{edo-3-3-2-acq-milestones-phases-2025}.
  \item[\textbf{JUON.}] \ac{juon} is a joint requirement identified by \ac{ccmd}, \ac{cjcs}, or \ac{vcjcs} for an ongoing contingency~\autocite{edo-3-3-2-acq-milestones-phases-2025}.
\end{description}

\subsection{DAS Structure: Milestones, Phases, Decision Points}
\begin{itemize}
  \item Milestones: Milestone~A initiates technology maturation and risk reduction; Milestone~B initiates engineering and manufacturing development; Milestone~C initiates production and deployment~\autocite{edo-3-3-2-acq-milestones-phases-2025}.
  \item Phases: \ac{msa}, \ac{tmrr}, \ac{emd}, \ac{pd.phase}, and \ac{os.phase}~\autocite{edo-3-3-2-acq-milestones-phases-2025}.
  \item Decision points: \ac{mdd}, \ac{cddv}, \ac{drfprd}, and \ac{frp}~\autocite{edo-3-3-2-acq-milestones-phases-2025}.
\end{itemize}

\begin{figure}[H]
  \centering
  \includegraphics[width=\linewidth]{MCA_lifecycle_view.png}
  \caption[Major Capability Acquisition Phases]{Major Capability Acquisition phases and decision points. \srcCite{edo-3-3-2-acq-milestones-phases-2025}.}
  \label{fig:mca_schedule_placeholder}
\end{figure}

\subsection{Milestone Review Information Requirements}
Information for milestone reviews is tailored by the \ac{mdauthority} and must align with DoWI~5000.85 requirements; it should cover program progress/risk, affordability, trade-offs, acquisition strategy updates, and exit criteria for the next phase~\autocite{edo-3-3-2-acq-milestones-phases-2025}.

\subsection{Entrance and Exit Criteria}
\begin{description}
  \item[\textbf{Entrance criteria.}] Minimum accomplishments required before entering a phase; documented in DoWI~5000.85 and required for all programs~\autocite{edo-3-3-2-acq-milestones-phases-2025}.
  \item[\textbf{Exit criteria.}] Program-specific technical, schedule, or management risk areas proposed by the \ac{pm} and approved by the \ac{mdauthority}; recorded in the \ac{adm} and must be specific and demonstrable~\autocite{edo-3-3-2-acq-milestones-phases-2025}.
  \item[\textbf{Risk linkage.}] Programs must meet exit criteria for the current phase and entrance criteria for the next phase before transition; the \ac{mdauthority} uses exit criteria to confirm risk reduction~\autocite{edo-3-3-2-acq-milestones-phases-2025}.
\end{description}

\subsection{Acquisition Program Baseline (APB) and Trade Space}
\begin{description}
  \item[\textbf{Purpose.}] The \ac{apb} documents cost, schedule, and performance thresholds/objectives and defines the trade space for the program~\autocite{edo-3-3-2-acq-milestones-phases-2025}.
  \item[\textbf{Thresholds.}] Minimum acceptable values; exceeding schedule, cost, or performance thresholds constitutes a deviation and requires \ac{mdauthority} involvement~\autocite{edo-3-3-2-acq-milestones-phases-2025}.
  \item[\textbf{Objectives.}] Desired values that represent meaningful improvements beyond thresholds and guide \ac{pm} trade decisions~\autocite{edo-3-3-2-acq-milestones-phases-2025}.
  \item[\textbf{Default thresholds.}] Unless specified, performance threshold equals objective, cost threshold equals objective plus 10\%, and schedule threshold equals objective plus 6~months~\autocite{edo-3-3-2-acq-milestones-phases-2025}.
\end{description}

\subsection{Program Deviations}
\begin{description}
  \item[\textbf{Deviation definition.}] A program deviation occurs when the current estimate exceeds \ac{apb} thresholds for cost, schedule, or performance; early indicators are set by the \ac{mdauthority} (e.g., 10\% schedule slip)~\autocite{edo-3-3-2-acq-milestones-phases-2025}.
  \item[\textbf{Nunn--McCurdy.}] Significant breaches are 15\% growth from the current baseline or 30\% from the original baseline; critical breaches are 25\% current or 50\% original, and are subject to termination review; significant breaches require congressional notification within 45 days after the Program Deviation Report~\autocite{edo-3-3-2-acq-milestones-phases-2025}.
\end{description}

\subsection{Program Deviation Actions}
\begin{itemize}
  \item Immediately notify the \ac{mdauthority} through appropriate management channels~\autocite{edo-3-3-2-acq-milestones-phases-2025}.
  \item Within 30~days, submit a Program Deviation Report describing the cause and corrective actions~\autocite{edo-3-3-2-acq-milestones-phases-2025}.
  \item Within 90~days, either return the program within \ac{apb} parameters or submit information to the Overarching \ac{ipt} for an \ac{mdauthority} decision on \ac{apb} revision~\autocite{edo-3-3-2-acq-milestones-phases-2025}.
\end{itemize}

\subsection{JCIDS Relationship to the DAS}
\ac{jcids} is the requirements generation process that identifies warfighter capability gaps, develops joint concepts, and produces a validated \ac{icd} through the Capability Based Assessment that feeds the acquisition framework~\autocite{edo-3-3-2-acq-milestones-phases-2025}.

\subsection{Materiel Development Decision (MDD)}
\begin{description}
  \item[\textbf{Purpose.}] Confirms a materiel solution is required and triggers the \ac{aoa}~\autocite{edo-3-3-2-acq-milestones-phases-2025}.
  \item[\textbf{MDA inputs.}] Validated \ac{icd}, evidence of technical foundation, \ac{aoa} study guidance and plan, and a Capstone Threat Assessment~\autocite{edo-3-3-2-acq-milestones-phases-2025}.
  \item[\textbf{MDA decisions.}] Designate the lead \ac{dod} Component and the phase entry point; entry may occur at any point that meets entrance criteria and statutory requirements~\autocite{edo-3-3-2-acq-milestones-phases-2025}.
\end{description}

\subsection{Materiel Solution Analysis (MSA)}
\begin{description}
  \item[\textbf{Purpose.}] Assess candidate materiel solutions and prepare for Milestone~A~\autocite{edo-3-3-2-acq-milestones-phases-2025}.
  \item[\textbf{Guided by.}] Validated \ac{icd}, \ac{aoa} study guidance, and \ac{aoa} plan~\autocite{edo-3-3-2-acq-milestones-phases-2025}.
  \item[\textbf{Major activities.}] Establish \ac{pm}/\ac{pmo}, conduct \ac{aoa}, develop initial Acquisition Strategy, \ac{temp}, \ac{sep}, \ac{lcsp}, and cybersecurity strategy; the user writes a draft \ac{cdd}~\autocite{edo-3-3-2-acq-milestones-phases-2025}.
  \item[\textbf{Minimum funding.}] Funding supports all phase activities and the Milestone~A decision~\autocite{edo-3-3-2-acq-milestones-phases-2025}.
  \item[\textbf{Exit.}] \ac{mdauthority} approves the materiel solution and Acquisition Strategy~\autocite{edo-3-3-2-acq-milestones-phases-2025}.
\end{description}

\subsection{Analysis of Alternatives (AoA)}
\begin{description}
  \item[\textbf{Focus.}] Materiel solution, mission effectiveness, key cost/capability trades, schedule, \ac{conops}, risk, producibility, and total life-cycle cost~\autocite{edo-3-3-2-acq-milestones-phases-2025}.
  \item[\textbf{Innovation/competition.}] Emphasizes competition and encourages industry to surface candidate solutions to sponsors and stakeholders~\autocite{edo-3-3-2-acq-milestones-phases-2025}.
  \item[\textbf{Hierarchy of alternatives.}] Commercial solutions first, then existing U.S./allied systems, international cooperative development, joint/government agency development, and finally \ac{dod}-unique development~\autocite{edo-3-3-2-acq-milestones-phases-2025}.
  \item[\textbf{Approval.}] \ac{cape} approves the \ac{aoa}; results inform the Acquisition Strategy and industry competition~\autocite{edo-3-3-2-acq-milestones-phases-2025}.
\end{description}

\subsection{Acquisition Strategy}
\begin{description}
  \item[\textbf{Requirement.}] The \ac{pm} prepares and the \ac{mdauthority} approves an Acquisition Strategy at Milestone~A~\autocite{edo-3-3-2-acq-milestones-phases-2025}.
  \item[\textbf{What it does.}] Tailors the acquisition framework, sets the master schedule, and provides the basis for plans such as \ac{temp}, contract awards, and \ac{sep}~\autocite{edo-3-3-2-acq-milestones-phases-2025}.
  \item[\textbf{Affordability.}] The strategy must maximize affordability from initiation through post-production support~\autocite{edo-3-3-2-acq-milestones-phases-2025}.
\end{description}

\subsection{International Acquisition and Exportability}
\begin{description}
  \item[\textbf{Drivers.}] U.S. strategy and DoWI~5000.02 require international cooperation, sales, and exportability to be considered in acquisition strategies~\autocite{edo-3-3-2-acq-milestones-phases-2025}.
  \item[\textbf{IA\&E objectives.}] Operational interoperability, economic cost-sharing, technical access, political alliance strength, and industrial base health; additional considerations include security risk and foreign dependency~\autocite{edo-3-3-2-acq-milestones-phases-2025}.
  \item[\textbf{Common forms.}] International cooperative programs, \ac{fms}, building partner capacity initiatives, and \ac{dcs} arrangements (do not confuse with program ownership and contract privity roles)~\autocite{edo-3-3-2-acq-milestones-phases-2025}.
\end{description}

\subsection{Evolutionary v.s.\ Single-Step Acquisition}
\begin{description}
  \item[\textbf{Evolutionary.}] Time-phased increments that define the current increment and subsequent evolutions to meet full capability needs~\autocite{edo-3-3-2-acq-milestones-phases-2025}.
  \item[\textbf{Single-step.}] Full capability delivered in one step before Milestone~C; major modifications of sufficient cost/complexity are managed as separate efforts~\autocite{edo-3-3-2-acq-milestones-phases-2025}.
  \item[\textbf{Preference.}] Evolutionary acquisition is preferred when rapidly fielding mature technology~\autocite{edo-3-3-2-acq-milestones-phases-2025}.
\end{description}

\subsection{Key Performance Parameters}
\begin{description}
  \item[\textbf{Definition.}] \acp{kpp} are system attributes critical to delivering effective capability and are documented in the \ac{cdd} and \ac{apb}~\autocite{edo-3-3-2-acq-milestones-phases-2025}.
  \item[\textbf{Validation.}] \ac{jroc} validates \acp{kpp} for major defense programs; the Joint Capabilities Board or \ac{dod} Components validate for others~\autocite{edo-3-3-2-acq-milestones-phases-2025}.
  \item[\textbf{Mandatory.}] Force Protection, System Survivability, Sustainment, and Energy; Net-Ready remains required but can be treated as a \ac{kpp} or performance attribute~\autocite{edo-3-3-2-acq-milestones-phases-2025}.
\end{description}

\subsection{Milestone A}
\begin{description}
  \item[\textbf{PM inputs.}] Proposed materiel solution, Acquisition Strategy, \ac{sep}/\ac{temp}, \ac{cdd}/\ac{conops}, \ac{lcsp}/\ac{ice}, risk assessment, and should-cost targets~\autocite{edo-3-3-2-acq-milestones-phases-2025}.
  \item[\textbf{Component inputs.}] Affordability analysis, cost estimate, and evidence the program is fully funded in the \ac{fydp}~\autocite{edo-3-3-2-acq-milestones-phases-2025}.
  \item[\textbf{MDA decisions.}] Approves materiel solution, authorizes \ac{tmrr} \ac{rfp} release, and sets \ac{tmrr} exit criteria in the \ac{adm}~\autocite{edo-3-3-2-acq-milestones-phases-2025}.
  \item[\textbf{Certification.}] 10~U.S.C.\ 2366a requires \ac{icd} alignment, \ac{aoa}-validated solution, cost estimate with CAPE concurrence, \ac{jroc} priority, and \ac{dia}-validated \ac{volt} report~\autocite{edo-3-3-2-acq-milestones-phases-2025}.
\end{description}

\subsection{Technology Maturation and Risk Reduction (TMRR)}
\begin{description}
  \item[\textbf{Purpose.}] Reduce technology, engineering, integration, and life-cycle cost risks; demonstrate critical technologies; complete preliminary design~\autocite{edo-3-3-2-acq-milestones-phases-2025}.
  \item[\textbf{Guided by.}] \ac{icd}, Acquisition Strategy, draft \ac{cdd}, and \ac{sep}~\autocite{edo-3-3-2-acq-milestones-phases-2025}.
  \item[\textbf{Basis for entry.}] \ac{mdauthority}-approved materiel solution and Acquisition Strategy~\autocite{edo-3-3-2-acq-milestones-phases-2025}.
  \item[\textbf{Activities.}] Competitive prototyping, \ac{pdr}, \ac{cddv}, sustainment planning, \ac{drfprd}, \ac{tra}, and \ac{ramc} report~\autocite{edo-3-3-2-acq-milestones-phases-2025}.
  \item[\textbf{Exit criteria.}] Affordable increment identified, technology demonstrated in relevant environment (\ac{trl}~6), \ac{pdr} completed before Milestone~B (unless waived), validated requirements, full \ac{fydp} funding, and affordability compliance for production/sustainment~\autocite{edo-3-3-2-acq-milestones-phases-2025}.
\end{description}

\subsection{Capability Development Document Validation}
\begin{description}
  \item[\textbf{Authority.}] The validation authority (\ac{jroc} for \acp{mdap}) validates the \ac{cdd}~\autocite{edo-3-3-2-acq-milestones-phases-2025}.
  \item[\textbf{Focus.}] Confirms cost-performance trades, affordability, testability, and alignment with user priorities; \acp{kpp}/\acp{ksa} guide \ac{pdr} and \ac{drfprd}~\autocite{edo-3-3-2-acq-milestones-phases-2025}.
\end{description}

\subsection{Configuration Steering Boards}
\begin{description}
  \item[\textbf{Authority.}] Chaired by the \ac{cae} with \ac{usdas} and Joint Staff participation; formed after \ac{cddv}~\autocite{edo-3-3-2-acq-milestones-phases-2025}.
  \item[\textbf{Responsibilities.}] Review requirement changes for \ac{acat}~I/IA programs, propose adjustments to meet affordability, and evaluate cost/schedule impacts; changes require identified funds and mitigated schedule impacts~\autocite{edo-3-3-2-acq-milestones-phases-2025}.
  \item[\textbf{Cadence.}] Meets at least annually and recommends de-scope options to the \ac{mdauthority}~\autocite{edo-3-3-2-acq-milestones-phases-2025}.
\end{description}

\subsection{Development RFP Release Decision}
The \ac{drfprd} authorizes release of the \ac{rfp} for \ac{emd} (often with \ac{lrip} options), confirms affordability/executability, and is the last point for major changes without program disruption; the \ac{mdauthority} determines preliminary \ac{lrip} quantities~\autocite{edo-3-3-2-acq-milestones-phases-2025}.

\subsection{Milestone B}
\begin{description}
  \item[\textbf{Purpose.}] Development decision authorizing entry into \ac{emd} and formal program initiation~\autocite{edo-3-3-2-acq-milestones-phases-2025}.
  \item[\textbf{MDA approvals.}] \ac{emd} entry, Acquisition Strategy, \ac{apb}, \ac{lrip} quantities, contract type, exit criteria, and Milestone~B certification/determination with \ac{adm} issuance~\autocite{edo-3-3-2-acq-milestones-phases-2025}.
  \item[\textbf{Shipbuilding.}] Milestone~B triggers \ac{scn} funding for shipbuilding programs~\autocite{edo-3-3-2-acq-milestones-phases-2025}.
\end{description}

\subsection{Engineering and Manufacturing Development}
\begin{description}
  \item[\textbf{Purpose.}] Develop, build, and test the product to verify all requirements and support production decisions~\autocite{edo-3-3-2-acq-milestones-phases-2025}.
  \item[\textbf{Guided by.}] Acquisition Strategy, \ac{cdd}, \ac{temp}, and \ac{sep}~\autocite{edo-3-3-2-acq-milestones-phases-2025}.
  \item[\textbf{Activities.}] Complete design, retire risks, build/test prototypes, establish product baselines, conduct \ac{cdr}, and perform operational assessments~\autocite{edo-3-3-2-acq-milestones-phases-2025}.
  \item[\textbf{Exit.}] Design stability, requirements met via \ac{dt} and early \ac{optest}, and all \ac{emd} exit criteria met before Milestone~C~\autocite{edo-3-3-2-acq-milestones-phases-2025}.
\end{description}

\subsection{Milestone C}
At Milestone~C the \ac{mdauthority} reviews cost estimates, validated life-cycle threat reporting, environmental issues, design stability, \ac{oa} results, software maturity, manufacturing risk, validated \ac{cdd}, supportability, interoperability, affordability, and phased production ramp-up before committing to production~\autocite{edo-3-3-2-acq-milestones-phases-2025}.

\subsection{Production and Deployment}
\begin{description}
  \item[\textbf{Purpose.}] Produce and deliver requirements-compliant systems~\autocite{edo-3-3-2-acq-milestones-phases-2025}.
  \item[\textbf{Guided by.}] Acquisition Strategy, \ac{temp}, \ac{cdd}, \ac{sep}, and \ac{lcsp}~\autocite{edo-3-3-2-acq-milestones-phases-2025}.
  \item[\textbf{Activities.}] \ac{lrip}, \ac{optest} in realistic threat conditions, Beyond \ac{lrip} report, \ac{frpdr}, fielding, \ac{ioc}, and \ac{foc}~\autocite{edo-3-3-2-acq-milestones-phases-2025}.
  \item[\textbf{Exit.}] \ac{ioc} achieved and deployment in full swing~\autocite{edo-3-3-2-acq-milestones-phases-2025}.
\end{description}

\subsection{Operations and Support}
\begin{description}
  \item[\textbf{Purpose.}] Execute the support strategy, sustain readiness, and plan disposal~\autocite{edo-3-3-2-acq-milestones-phases-2025}.
  \item[\textbf{Guided by.}] Product Support Strategy and \ac{peshe}~\autocite{edo-3-3-2-acq-milestones-phases-2025}.
  \item[\textbf{Activities.}] Life-cycle sustainment resourcing, technical data/tooling planning, and disposal in accordance with safety, security, and environmental requirements; phase completes when assets are disposed~\autocite{edo-3-3-2-acq-milestones-phases-2025}.
\end{description}

\subsection{Ships Are Different}
\begin{itemize}
  \item The \ac{pm} is designated as early as possible, often before the \ac{mdd}~\autocite{edo-3-3-2-acq-milestones-phases-2025}.
  \item Initial Acquisition Strategy and \ac{temp} set the long-term context and gain early leadership alignment~\autocite{edo-3-3-2-acq-milestones-phases-2025}.
  \item Between Milestone~A and Milestone~B/C, \ac{rfp} and specification development is for the ship, not an engineering development model~\autocite{edo-3-3-2-acq-milestones-phases-2025}.
\end{itemize}

\subsection{Navy 2-Pass / 7-Gate Review}
\begin{description}
  \item[\textbf{Gates 1--4 (requirements).}] Led by the \ac{cno} or \ac{cmc}; starts before \ac{mdd} and leads to approved \ac{icd}, \ac{aoa} guidance, a selected \ac{aoa} alternative, approved \ac{cdd}, \ac{conops}, and a System Design Specification development plan~\autocite{edo-3-3-2-acq-milestones-phases-2025}.
  \item[\textbf{Gates 5--7 (acquisition/sustainment).}] Led by \ac{asnrdanda}; starts after Gate~4 and continues through Milestone~B; approves System Design Specification, \ac{rfp} release, and production readiness (including \ac{evms} and \ac{ibr} checks), with Gate~6 pre/post Milestone~C and \ac{frpdr} follow-on~\autocite{edo-3-3-2-acq-milestones-phases-2025}.
\end{description}

\subsection{Scenario Drills (Entry-Point Decisions)}
Use the following vignettes to select the most appropriate entry point into the acquisition framework~\autocite{edo-3-3-2-acq-milestones-phases-2025}.
\begin{enumerate}
  \item \textbf{Situation 1.} Joint mid-course interceptor capability; validated \ac{icd}; multiple materiel approaches; \ac{aoa} guidance and plan approved; MDA wants a lead \ac{dod} Component and an evolutionary strategy. \emph{Question:} best entry point?
  \item \textbf{Situation 2.} Validated \ac{cdd} for a new side-arm with proven commercial options; field evaluations meet thresholds; full procurement funding. \emph{Question:} best entry point?
  \item \textbf{Situation 3.} Validated \ac{icd}; \ac{aoa} and Acquisition Strategy complete; promising but immature technology; significant performance risk; draft \ac{cdd}. \emph{Question:} best entry point?
  \item \textbf{Situation 4.} Mature lab eye-protection technology; validated \ac{icd}/\ac{cdd}; full funding offsets identified; technology not yet tested outside the lab and not integrated with helmet systems. \emph{Question:} best entry point?
\end{enumerate}

\subsection{Quick Review}
\begin{itemize}
  \item Which entry point requires a validated \ac{icd}, \ac{aoa} study guidance, and \ac{aoa} plan~\autocite{edo-3-3-2-acq-milestones-phases-2025}?
  \item Which milestone authorizes entry into \ac{pd.phase}~\autocite{edo-3-3-2-acq-milestones-phases-2025}?
  \item Which milestone requires all sources of program risk to be adequately mitigated~\autocite{edo-3-3-2-acq-milestones-phases-2025}?
  \item In \ac{msa}, which two artifacts are approved by the \ac{mdauthority}~\autocite{edo-3-3-2-acq-milestones-phases-2025}?
\end{itemize}

%====================
% End of file
%====================
\IfSubfilesClassLoaded{
  \RenewDocumentCommand{\entryname}{}{\textbf{\color{Modern} Acronym}}
  \RenewDocumentCommand{\descriptionname}{}{\textbf{\color{Modern} Definition}}
  \printnoidxglossary[
    type=\acronymtype,
    title=Acronyms,
    style=long-booktabs
  ]}{}
\end{document}
